\section{Prescribed charges on the full numerical lattice}
\label{sec:ambient}

We begin the proof of Theorem~\ref{thm:geometric-sector} by constructing
both prescribed charges with full numerical support.  This argument is
independent of the strong tilt inequalities.  Throughout the section,
$X$ is a smooth projective complex threefold with $K_X\simeq\OO_X$.

The construction takes place first on a product of projective spaces.
We choose an embedding that retains every numerical class of $X$, prove
a support estimate adapted to increasing volume on the product, and use
that estimate to deform its charge.  A fixed polynomial correction then
makes Grothendieck--Riemann--Roch restrict the deformed charge to the
chosen charge on $X$.  The same construction on smooth surface sections
will later identify the ordinary heart.
We use numerical Grothendieck groups in the Euler-pairing sense:
\[
 \Knum(X)=K_0(X)/\{v:\chi(v,u)=0\text{ for every }u\in K_0(X)\}.
\]
The left and right radicals coincide by Serre duality, so either side of
the Euler pairing may be used in the adjunction argument below.
All real vector spaces associated with these groups are finite-dimensional.
We use Bridgeland's stability conditions and locally finite slicings
\cite{Bridgeland2007}.  A stability condition has \emph{full numerical support} if, for a norm on
$\Knum(X)_\RR$, there is a constant $C$ such that
$\|[E]\|\le C|Z(E)|$ for every semistable object $E$.
The constant may depend on the stability condition.  This is equivalent to
the quadratic-form definition
\cite[Definition~1.1]{BayerDeformation2019}: one may take
$C^2|Z(v)|^2-\|v\|^2$, which is negative definite on $\ker Z$.

\subsection{An embedding which retains every numerical class}

The product-projective embedding and full-support construction here adapt
the approach of Giovenzana--Robotis--Rota--Zuliani
\cite{GiovenzanaRobotisRotaZuliani2026}.  The numerical injectivity argument
is given below; the quantitative product inputs and the prescribed-charge,
real-sector refinements are treated separately in the subsequent subsections.

\begin{lemma}\label{lem:full-numerical-embedding}
Let $X$ be a smooth projective threefold.  Rational Chern characters identify
its numerical Grothendieck group with
\begin{equation}\label{eq:full-numerical-lattice}
 \Knum(X)_\QQ\simeq
 \QQ\oplus N^1(X)_\QQ\oplus N_1(X)_\QQ\oplus\QQ.
\end{equation}
Products of numerical divisor classes span this space.  Given a real ample
class $H_*$, there are very ample line bundles $L_1,\ldots,L_q$ whose classes
form a basis of $N^1(X)_\RR$ and positive numbers $s_1,\ldots,s_q$ with
$H_*=\sum_hs_hL_h$.  For the product embedding
\begin{equation}\label{eq:product-embedding}
 i:X\hookrightarrow Y=\prod_{h=1}^q\mathbb P^{d_h},
 \qquad d_h=h^0(X,L_h)-1,
\end{equation}
the pushforward $i_*:\Knum(X)_\RR\to\Knum(Y)_\RR$ is injective.
\end{lemma}

\begin{proof}
The rational Chern character identifies the Grothendieck group with the
rational Chow group
\cite[Tags 0FDI, 0FEW, 0FB2 and 0FB5]{StacksProject}.  In Riemann--Roch the Euler pairing is the intersection
pairing after applying the invertible operations of dualizing a class and
multiplying it by $\td(X)$.  Its radical therefore gives precisely the
numerical quotient of the Chow group.  In dimensions zero and three this
quotient is $\QQ$; in the other two dimensions it is the divisor--curve
pairing in \eqref{eq:full-numerical-lattice}.

For a rational ample class $A$, the Hodge index theorem makes
$(D,D')\mapsto ADD'$ nondegenerate on $N^1(X)_\QQ$; see
\cite[Theorem~6.4 and Corollary~6.5]{GrebRossToma2016} for its
real-ample formulation.  The divisor--curve
pairing is nondegenerate by numerical equivalence, so multiplication by $A$
is an isomorphism $N^1(X)_\QQ\to N_1(X)_\QQ$.  Divisors, their products with
$A$, and $A^3$ consequently span every summand of
\eqref{eq:full-numerical-lattice}.

Choose an affine slice transverse to the ray of $H_*$ and a rational simplex
in its ample cone whose interior contains that ray.  Its $q=\rho(X)$
vertices are linearly independent.  Taking sufficiently large integral
multiples makes them very ample without changing their positive cone.  This
also explains the rank-one case, where a single very ample class suffices.

Pushforward is defined on numerical groups by adjunction with perfect
complexes.  If $i_*v=0$ numerically, the projection formula gives
\[
 \chi(i^*W,v)=\chi(W,i_*v)=0\qquad(W\in K_0(Y)).
\]
Chern characters of line bundles on $Y$ span its polynomial Chow ring, and
their restrictions span \eqref{eq:full-numerical-lattice}.  Nondegeneracy of
the Euler pairing gives $v=0$.  The spanning argument here uses dimension
three; it makes no analogous claim in arbitrary dimension.
\end{proof}

\subsection{Scaled support on elliptic products}

The next estimate is the quantitative reason that the ambient charge can
be corrected uniformly at large volume.  Its coordinates rescale a
codimension-$j$ character by the expected power $R^{n-j}$.  In these
coordinates a fixed positive-codimension correction of the exponential
charge will tend to zero as $R\to\infty$.  A support bound independent
of $R$ then turns this coordinate estimate into a bound on every
semistable object.

We first establish the support bound on an elliptic product.  The product
construction and the uniqueness needed to compare its different orders
of induction are external inputs.  The estimate below makes their volume
scaling explicit.

\begin{lemma}[Scaled support on an elliptic product]
\label{lem:elliptic-scaled-support}
Let $C$ be an elliptic curve, let $n\ge1$, and fix positive rational
numbers $\rho_1,\ldots,\rho_n$.  For rational $R\ge1$, put
$w_r=R\rho_r$ and let $H_r$ be the point-divisor class from the $r$th
factor of $C^n$.  The product stability condition with charge
\[
 Z_{\mathbf w}(u)=-\int_{C^n}e^{-i\sum_rw_rH_r}\ch(u)
\]
has support on the coordinate lattice
\[
 x_I(u)=H_I\ch_{n-|I|}(u),\qquad
 H_I=\prod_{r\in I}H_r,\qquad I\subset\{1,\ldots,n\}.
\]
More precisely, set $y_I=(\prod_{r\in I}w_r)x_I$.  In these coordinates,
\[
 Z_{\mathbf w}(u)=\widehat Z_n(y(u)),\qquad
 \widehat Z_n(y)=-\sum_I(-i)^{|I|}y_I.
\]
For a fixed Euclidean norm there is $C_n>0$, independent of $R$, such that
\[
 \|y(u)\|\le C_n|Z_{\mathbf w}(u)|
 \qquad(u\text{ the class of a semistable object}).
\]
The constants can be chosen uniformly when the positive ratios $\rho_r$
vary in a compact set.  Point sheaves are stable of phase $1$.
\end{lemma}

\begin{proof}
We use the product-with-a-curve construction of
\cite[Lemma~5.7 and Remark~5.8]{LiuProduct2021}, the positive-genus
product construction of \cite[Theorem~4.5]{ProductsCurves2025}, and
the support induction of \cite[Lemma~6.6]{LiEtAl2026}.
For $n=1$, the usual slope stability on $C$ has charge
$-x_{\varnothing}+iw_1x_{\{1\}}$ and the assertion is immediate.

Suppose that the scaled support bound has been proved in dimension
$n-1$.  Omitting factor $r$ splits the coordinates into $x',x''$,
where $x'$ consists of those containing $r$, with that index removed.
The product construction has charge
\[
 Z_{n-1}(x'')-iw_rZ_{n-1}(x').
\]
It initially gives support with respect to its own quotient lattice.
Every choice of omitted factor gives the same exponential charge and
the same phase $1$ for point sheaves.  The uniqueness theorem
\cite[Theorem~3.11]{ProductsCurves2025} therefore identifies the
slicings: that theorem permits the two support lattices to differ.
We may consequently test all the product support forms on the same
semistable objects, before proving support on the full coordinate
lattice.

Write $M_{\widehat r}$ for scaling the remaining coordinates by their
products of weights.  The scaled coordinates split as
\[
 y'=w_rM_{\widehat r}x',\qquad y''=M_{\widehat r}x''.
\]
The lower-dimensional charge in scaled coordinates is a fixed linear
map $\widehat Z$.  Increase its support constant $C_{n-1}$ if needed
and choose the support form
\begin{equation}\label{eq:elliptic-support-form}
 q_{n-1}(v)=C_{n-1}^2|\widehat Z(v)|^2-\|v\|^2.
\end{equation}
It is nonnegative on lower-dimensional semistable classes.  Crucially,
\[
 q_{n-1}(v)\le C_{n-1}^2|\widehat Z(v)|^2
 \quad\text{for every }v,
\]
not merely for semistable classes.  This is the norm form used in the
product support estimate.

For clarity, that estimate has the following coefficient condition.
Write $Z_{n-1}(x')=\alpha+i\beta$ and
$Z_{n-1}(x'')=\gamma+i\delta$.  With equal positive product parameters
$s=t=w_r$, the form
\[
 \beta\gamma-\alpha\delta
       +\eta_r q_{n-1}(M_{\widehat r}x')
\]
is nonnegative on product-semistable classes when
$0\le\eta_r\le w_r/C_{n-1}^2$
\cite[Lemma~5.7]{LiuProduct2021}; this is the substitution in
\cite[Lemma~6.6]{LiEtAl2026}.
Choose $0<\lambda_r<C_{n-1}^{-2}$, independently of $R$, and put
$\eta_r=\lambda_rw_r$.  If
\[
 \widehat Z(y')=\widehat\alpha+i\widehat\beta,
 \qquad \widehat Z(y'')=\widehat\gamma+i\widehat\delta,
\]
then multiplication of the product form by $w_r$ gives
\begin{equation}\label{eq:weighted-product-form}
 q_r(y)=\widehat\beta\widehat\gamma
       -\widehat\alpha\widehat\delta
       +\lambda_rq_{n-1}(y').
\end{equation}
All powers of $R$ have disappeared.  This form is nonnegative on the
common semistable objects just constructed.

The total charge in these coordinates is
$(\widehat\gamma+\widehat\beta)
+i(\widehat\delta-\widehat\alpha)$.  On its kernel,
\[
 q_r(y)=-|\widehat Z(y')|^2+\lambda_rq_{n-1}(y')
       =(\lambda_rC_{n-1}^2-1)|\widehat Z(y')|^2
           -\lambda_r\|y'\|^2.
\]
It is nonpositive there, and vanishes only when $y'=0$; the kernel
condition then also gives $\widehat Z(y'')=0$.
Sum \eqref{eq:weighted-product-form} over all omitted factors.  If the
sum vanishes on the total charge kernel, then $y'=0$ for every factor.
This kills every coordinate with $I\ne\varnothing$, and the charge
kills the remaining top-degree coordinate.  The summed form is thus
negative definite on the charge kernel and nonnegative on semistables.
Compactness on the unit sphere of its nonnegative cone gives
$\|y(u)\|\le C_n|Z_{\mathbf w}(u)|$ for every semistable class $u$.
The forms and all choices were made in the scaled coordinates.  In
particular they are independent of $R$ and can be kept uniform on
compact sets of positive ratios.  This completes the support induction.
The geometricity assertion is part of the positive-genus product
construction; see also \cite[Theorem~4.1(1)]{LiEtAl2026}.
\end{proof}

\subsection{Transfer to projective blocks}

For a slicing $\mathcal R$, denote the largest and smallest
Harder--Narasimhan phases of a nonzero object by $\phi^+_{\mathcal R}$ and
$\phi^-_{\mathcal R}$.  We use the distance
\[
 d(\mathcal R,\mathcal R')=
 \sup_{E\ne0}\max\{
 |\phi^+_{\mathcal R}(E)-\phi^+_{\mathcal R'}(E)|,
 |\phi^-_{\mathcal R}(E)-\phi^-_{\mathcal R'}(E)|\}.
\]
For a line bundle $L$ and an integer $k$, we say that $\mathcal R$ has
the \emph{Bayer property} for $(L,k)$ if $F\otimes L[k]$ has all phases
at least $\phi$ whenever $F\in\mathcal R(\phi)$.
We will use both a support bound and a small phase change under each
fixed line-bundle twist.  They are inherited from an elliptic product
through finite quotients.

\begin{proposition}\label{prop:weighted-ambient}
Let $Y=\prod_{h=1}^q\mathbb P^{d_h}$, put
$\xi_h=c_1(\OO_h(1))$, and fix $s_h>0$.  There is a continuous family of
numerical stability conditions $(U_R,\mathcal R_R)$ for $R\ge1$ with
\[
 U_R(u)=-\int_Ye^{-iR\sum_hs_h\xi_h}\ch(u).
\]
All point sheaves are stable of phase $1$.  Write
$\ch(u)=\sum_{\mathbf j}u_{\mathbf j}\xi^{\mathbf j}$, where
$0\le j_h\le d_h$, and set
\begin{equation}\label{eq:ambient-scaled-coordinates}
 (M_Ru)_{\mathbf j}=
 u_{\mathbf j}\prod_h(Rs_h)^{d_h-j_h}.
\end{equation}
For a fixed Euclidean norm there is $D>0$, independent of $R$, such that
\begin{equation}\label{eq:ambient-support}
 \|M_Ru\|\le D|U_R(u)|\qquad
 (u\text{ the class of an }\mathcal R_R\text{-semistable object}).
\end{equation}
For every $\mathbf e\in\ZZ^q$,
\begin{equation}\label{eq:ambient-twist-distance}
 d(\mathcal R_R,\mathcal R_R\otimes\OO_Y(\mathbf e))
 \le\sum_h\frac{d_h|e_h|}{\pi Rs_h}.
\end{equation}
\end{proposition}

\begin{proof}
Cheng proves the two-block construction and phase estimate in
\cite[Proposition~4.3 and Remark~4.4]{Cheng2026}.  We apply the same
quotient construction block by block, using
Lemma~\ref{lem:elliptic-scaled-support} to control all coordinates
uniformly in the volume.

First take $R$ and the positive weights $s_h$ rational.  Put
$n=\sum_hd_h$ and give
each of the $d_h$ elliptic factors in block $h$ the weight
$w_r=2Rs_h$.  The lemma supplies the starting stability condition and
its scaled support bound.  Only the coordinate lattice of the elliptic
power is used here; full numerical support will follow on $Y$ from the
explicit pullback calculation below.

First apply \cite[Proposition~5.1]{LiEtAl2026} to the stability condition
just constructed on the coordinate lattice.  It supplies the descent
by $(\ZZ/2)^n$ to $(\mathbb P^1)^n$ and the Bayer property with $k=0$
for every effective line bundle $\OO(a_1,\ldots,a_n)$, with all $a_r\ge0$.
We then take the symmetric quotients one block at a time.

Here is the filtration used at a symmetric step.  For
$Q=(\mathbb P^1)^d\to\mathbb P^d$, there is a finite filtration of
$\OO_{Q\times_{\mathbb P^d}Q}$ whose successive quotients are
\[
 p_2^*\mathcal L_j^{-1}\otimes\OO_{\Gamma(g_j)},
 \qquad g_j\in S_d,\qquad
 \mathcal L_j=\OO(a_{j1},\ldots,a_{jd}),\quad a_{jr}\ge0,
\]
where $p_2$ is the second projection and $\Gamma(g_j)$ is the graph
of the permutation $g_j$
\cite[Definition~3.19 and Example~3.20(3)]{LiEtAl2026}.
At a block stage the map is $f:Q\times M\to\mathbb P^d\times M$,
where $M$ is the product of the other blocks, and its fiber square is
$(Q\times_{\mathbb P^d}Q)\times M$.
Flat pullback to this product preserves the filtration and its exact
sequences.  Its graph automorphisms are $g_j\times\operatorname{id}_M$,
and its line bundles remain pulled back from $Q$.

The finite-flat descent criterion \cite[Proposition~3.21]{LiEtAl2026}
therefore requires precisely invariance under these permutations and
the Bayer properties with $k=0$ for these effective line bundles.
Indeed, for a phase-$\phi$ object $F$, the factors of $f^*f_*F$ are
$\mathcal L_j^{-1}\otimes g_{j*}F$; the two properties put their largest
phases at most $\phi$.  On the elliptic
product, signs and permutations within each equal-weight block preserve
the charge and point phases, so the uniqueness theorem used above gives
their invariance.  Block permutations normalize the full sign group;
their invariance descends through the completed involution quotient.
For subsequent block quotients, the required properties pass to the
remaining blocks as follows
\cite[Lemma~3.9 and Remark~3.14(3)]{LiEtAl2026}.
At each quotient $f:V\to V'$, pullback computes extremal phases for the
induced slicing $\mathcal S'=f_\#\mathcal S$. If a line bundle $L$ comes
from an unquotiented block, then
\[
 \phi^-_{\mathcal S'}(F\otimes L[k])
 =\phi^-_{\mathcal S}(f^*F\otimes f^*L[k]).
\]
The Bayer property for $f^*L$ therefore descends to $L$. An automorphism
of another block commutes with $f$, so its invariance descends by the
same pullback description. Thus every remaining quotient has
the required Bayer properties and invariance, for any finite number of
blocks.

For completeness, let $g$ be the resulting map from the elliptic power to
$Y$.  Pullback satisfies
\[
 g^*\xi_h^{j_h}=2^{j_h}j_h!
       \sum_{|J_h|=j_h}H_{J_h},
 \qquad \deg g=2^n\prod_hd_h!.
\]
If $|I\cap\text{block }h|=d_h-j_h$, then the scaled $I$-coordinate of
$g^*u$ is
\begin{equation}\label{eq:block-coordinate-comparison}
 \Bigl(\prod_{r\in I}w_r\Bigr)
 H_I\ch_{n-|I|}(g^*u)
 =2^n\Bigl(\prod_hj_h!\Bigr)(M_Ru)_{\mathbf j}.
\end{equation}
The fixed multiplicities are $\prod_h\binom{d_h}{j_h}$.  Pullback is thus
injective on the complete Chern-character lattice of $Y$, with norm
comparison constants independent of $R$.  Semistability pulls back, and
$Z_{\mathbf w}(g^*u)=\deg(g)U_R(u)$.  This proves \eqref{eq:ambient-support}.

For the phase estimate, vary the real tensor parameter on each elliptic
factor successively.  The one-factor argument in the proof of
\cite[Theorem~2.4]{ChengFeyzbakhsh2026}, used in
\cite[Remark~4.4]{Cheng2026}, gives its absolute variation divided by
$\pi w_r$.  A twist $e_h\xi_h$ pulls back with coefficient $2e_h$ on the
$d_h$ factors in block $h$; its factor $2$ cancels the factor $2$ in $w_r$.
Adding the variations gives \eqref{eq:ambient-twist-distance}.
Finally the deformation and multiplication-isogeny argument in the proof of
\cite[Proposition~4.3]{Cheng2026} extends the rational-weight construction
continuously to positive real weights.  It applies factor by factor: the
scaled support forms just constructed remain fixed on compact positive
weight sets, the exponential charges vary continuously, and uniqueness on
overlapping deformation neighborhoods identifies the lifts.  The closed
support inequalities and the phase bounds persist.  This also extends $R$
to all real $R\ge1$.  Point stability in this extension follows separately
from \cite[Theorem~4.1(1)]{LiEtAl2026} on the elliptic power and preservation
of geometricity under finite pushforward
\cite[Remark~3.11]{LiEtAl2026}.  Concretely, for the finite faithfully flat
quotient $g$, the sheaf $g^*\OO_y$ has a finite filtration by point sheaves,
all stable of phase $1$.  A phase-one subobject and quotient of $\OO_y$
pull back to objects whose Jordan--H\"older factors are among these points,
so both pullbacks are coherent finite-length sheaves.  Faithful flatness
descends this cohomological concentration.  The original subobject and
quotient are therefore sheaves, and simplicity of $\OO_y$ rules out a
proper nonzero subobject.  Continuity and the charge value $-1$ normalize
the point phase to $1$.
\end{proof}

\subsection{Geometric slicings and the restriction criterion}

We now prepare the restriction step.  The first consequence of point
stability below gives the cohomological bounds needed for a truncated
resolution.  The second will keep points stable during deformation in
all ambient charges.  Both will also be used to identify the induced
heart.  Write
$\mathcal R(I)$ for the extension-closed category generated by semistable
objects with phases in an interval $I$.

\begin{lemma}[Cohomological bounds from points]\label{lem:geometric-range}
Let $V$ be a smooth projective variety of dimension $n\ge2$, and let
$\mathcal R$ be a locally finite slicing in which every point sheaf is
stable of phase $1$.  Then
\begin{align}
 \mathcal R(0,1]&\subset D^{[-n+1,0]}(V),
 &\Coh(V)&\subset\mathcal R(1-n,1],\label{eq:geometric-range}\\
 D^{\le-n}(V)&\subset\mathcal R(>0),
 &\mathcal R(>0)&\subset D^{\le0}(V).\label{eq:geometric-aisles}
\end{align}
The sheaf $\mathcal H^{-n+1}(E)$ is torsion-free for
$E\in\mathcal R(0,1]$.  Moreover, $\mathcal H^0(E)$ is zero-dimensional
for $E\in\mathcal R(1)$.
\end{lemma}

\begin{proof}
This is the point-object argument of \cite[Lemma~10.1]{BridgelandK3}; we
include it to specify the dimension and boundary statements being used.
Let $E$ be stable with phase $0<\phi\le1$ and not a point of phase $1$.
Phase vanishing gives $\Hom(E,\OO_x[j])=0$ for $j<0$.
For $j\ge n$, Serre duality identifies this space with
$\Hom(\OO_x,E[n-j])^*$, which is zero by phases.  At $j=n,\phi=1$, use
nonisomorphism of the two stable objects.  The canonical bundle causes no
change, since its restriction to a point is trivial.  A minimal free
complex at each closed point consequently has nonzero terms only in
degrees $[-n+1,0]$.  Its bottom cohomology is a subsheaf of a locally free
sheaf and hence is torsion-free.  At phase $1$ we also have
$\Hom(E,\OO_x)=0$ for every $x$, so $\mathcal H^0(E)=0$.
Point sheaves themselves satisfy the stated cohomological bounds.  Finite
Jordan--Hölder filtrations and extensions give these assertions for the whole
heart; in phase $1$ its zeroth cohomology is an extension of sheaves supported
at finitely many points.

A coherent sheaf cannot receive a nonzero morphism from an object whose
phase is greater than $1$, by the cohomological bound just proved and a
shift.  It cannot map nontrivially to an object of phase at most $1-n$:
such an object is in $D^{\ge1}$, including the endpoint by the phase-one
observation.  Applied to its extreme Harder--Narasimhan factors, these
vanishings prove $\Coh(V)\subset\mathcal R(1-n,1]$.  Shifting this
inclusion and the heart bound proves \eqref{eq:geometric-aisles}.
\end{proof}

\begin{lemma}[Stability from an open set of charges]\label{lem:primitive-point}
Let a stability condition have support on a finite-rank lattice $\Lambda$.
Suppose that an object $E$ has primitive class in $\Lambda$ and remains
semistable throughout a neighborhood which is open in all central charges
$\Hom(\Lambda,\CC)$.  Then $E$ is stable.
\end{lemma}

\begin{proof}
If $E$ is strictly semistable, take a first stable subobject $G$ in a finite
Jordan--Hölder filtration.  Its class cannot be real proportional to $[E]$.
Indeed their aligned charges would make $[G]=c[E]$ with $0<c<1$, whereas
primitivity makes an integral class on this line an integral multiple of
$[E]$.  Choose a phase-window heart of length one centered at their common
phase.  The triangle $G\to E\to E/G$ remains a short exact sequence in
this heart for small deformations: all three objects stay strictly inside
its phase window.  An independent small perturbation of the two charge
values makes the phase of $G$ greater than that of $E$.  This contradicts
the assumed semistability of $E$ in the entire neighborhood.
\end{proof}

Here is the precise restriction input.  Let $f:V\to Y$ be a finite map.
For the truncated cofiltrations used below, the data consist of a finite
tower in $\Db(Y)$
\[
 E_0=f_*\OO_V\longrightarrow E_1\longrightarrow\cdots
      \longrightarrow E_N
\]
and distinguished triangles
\[
 P_j\longrightarrow E_j\longrightarrow E_{j+1}
      \longrightarrow P_j[1],\qquad 0\le j<N.
\]
Each $P_j$ is a finite direct sum of specified line bundles with specified
shifts $L_i[k_i]$, and the remainder satisfies $E_N\in D^{\le-N}(Y)$.
Thus the successive cones of the tower are $P_j[1]$; equivalently,
$f_*\OO_V$ is built by extensions from the $P_j$ and the remainder.
For a truncated resolution, $P_j$ is its $j$th sheaf term shifted by $[j]$
and $E_N$ is the last kernel shifted by $[N]$.
The finite-map restriction
criterion \cite[Proposition~3.24]{LiEtAl2026}, extending
\cite[Corollary~2.2.2]{Polishchuk2007}, applies to $f$ with this
cofiltration.  It requires
\begin{equation}\label{eq:restriction-hypotheses}
 \begin{gathered}
 \text{the Bayer properties for every }(L_i,k_i),\\
 D^{\le n_1}(Y)\subset\mathcal R(>0)\subset D^{\le n_2}(Y),
 \qquad N\ge n_2-n_1+1.
 \end{gathered}
\end{equation}
An untruncated cofiltration has no last numerical requirement.  All tensor
products in this criterion are derived.  The conclusion is a slicing
\begin{equation}\label{eq:induced-slicing}
 f^\sharp\mathcal R(\phi)=
       \{F\in\Db(V):f_*F\in\mathcal R(\phi)\},
\end{equation}
with charge $U\circ f_*$ and support on the ambient lattice through $f_*$.
In particular it provides the required Harder--Narasimhan filtrations.
No surjectivity assumption is imposed on $f$.  Smoothness makes the bounded
finite-Tor-dimension formulation applicable here.  We shall verify the
cofiltration and every bound in \eqref{eq:restriction-hypotheses} before
using the criterion.

\subsection{Correcting the charge before restriction}

Fix real classes $B_*,H_*$ with $H_*$ ample, and choose the embedding in
Lemma~\ref{lem:full-numerical-embedding}.  We write $L_h$ also for its first
Chern class, so $i^*\xi_h=L_h$ and $H_*=\sum_hs_hL_h$ with $s_h>0$.
Consider separately the factors
\begin{equation}\label{eq:two-character-factors}
 \Theta=1\qquad\text{and}\qquad
 \Theta=\sqrt{\td(X)}=1+\frac{c_2(X)}{24}.
\end{equation}
The assumption on the canonical bundle gives
$\td(X)=1+c_2(X)/12$ in the degrees relevant on $X$.  By
Lemma~\ref{lem:full-numerical-embedding}, choose a degree-two polynomial
$\delta$ in the $\xi_h$ with $i^*\delta=c_2(X)$ numerically.  A suitable
polynomial on $Y$ with constant term one is
\[
 \gamma=\begin{cases}
 \td(Y)(1-\delta/12),&\Theta=1,\\
 \td(Y)(1-\delta/24),&\Theta=\sqrt{\td(X)}.
 \end{cases}
\]
Products are taken in the finite-dimensional numerical Chow ring.  Thus
\begin{equation}\label{eq:gamma-lift}
 i^*\bigl(\gamma/\td(Y)\bigr)\td(X)=\Theta.
\end{equation}
In particular the two choices remove different Todd contributions.

Write $B=\sum_hb_hL_h$ and $H=\sum_h\ell_hL_h$.  Deform the ambient charge
to
\begin{equation}\label{eq:deformed-ambient-charge}
 U_t^{B,H}(u)=-\int_Y
 e^{-\sum_hb_h\xi_h-it\sum_h\ell_h\xi_h}\ch(u)\gamma.
\end{equation}
Grothendieck--Riemann--Roch and \eqref{eq:gamma-lift} give the exact identity
\begin{equation}\label{eq:exact-restricted-charge}
 U_t^{B,H}(i_*E)=-\int_Xe^{-B-itH}\ch(E)\Theta.
\end{equation}
Numerical equality in \eqref{eq:gamma-lift} suffices because only
intersection degrees occur.  For $\Theta=1$ this is exactly $Z_{B,tH}$;
for the other choice it is exactly $Z^{\mathrm{LV}}_{B,tH}$.

\begin{lemma}\label{lem:uniform-ambient-deformation}
For every $0<\epsilon<1/8$, there are open neighborhoods $V\ni B_*$ and
$W\ni H_*$ in the real divisor and ample spaces, and $T\ge1$, such that for all
$B\in V$, $H\in W$, and $t\ge T$, the charge
\eqref{eq:deformed-ambient-charge} lifts to a stability condition with
slicing at distance less than $\epsilon$ from
$\mathcal R_t$.  The construction is available in an open neighborhood of
that charge in the full space $\Hom(\Knum(Y),\CC)$.
\end{lemma}

\begin{proof}
Take $V$ bounded in the $b_h$ coordinates, and restrict $W$ by
$|\ell_h/s_h-1|<\eta$.  In the scaled coordinates
\eqref{eq:ambient-scaled-coordinates}, expansion of the two exponentials
gives
\begin{equation}\label{eq:uniform-deformation-estimate}
 \bigl\|(U_t^{B,H}-U_t)\circ M_t^{-1}\bigr\|
       \le C_1\eta+C_2t^{-1}.
\end{equation}
To see the uniformity, the coefficient of $u_{\mathbf j}$ in the
uncorrected exponential has degree $n-|\mathbf j|$ in $t$.  Varying the
positive coefficients $s_h$ to $\ell_h$ changes its ratio to the weight of
$M_t$ by $O(\eta)$.  A codimension-$d>0$ term of
$e^{-\sum b_h\xi_h}\gamma-1$ lowers that degree by $d$ and hence contributes
$O(t^{-d})$.  There are finitely many monomials, and their coefficients are
bounded on the chosen neighborhoods.  This proves
\eqref{eq:uniform-deformation-estimate} independently of the object.

Combine it with \eqref{eq:ambient-support}.  Choose $\eta$ first and then
$T$ so that
\[
 D(C_1\eta+C_2/T)<\sin(\pi\epsilon).
\]
To check the support hypothesis of the deformation theorem explicitly, put
\[
 Q_t(u)=D^2|U_t(u)|^2-\|M_tu\|^2.
\]
It is nonnegative on the original semistable classes.  If $U'(u)=0$ and
$\|(U'-U_t)M_t^{-1}\|\le\delta<1/D$, then
\[
 Q_t(u)\le(D^2\delta^2-1)\|M_tu\|^2<0\qquad(u\ne0).
\]
Thus $Q_t$ is negative definite on every charge kernel along the straight
deformation, and throughout the indicated open charge region.
The support-property deformation theorem
\cite[Theorem~1.2]{BayerDeformation2019} lifts the straight path while
preserving $Q_t$.  On its nonnegative cone,
$\|M_tu\|\le D|U_t(u)|$, so the relative charge change is less than
$\sin(\pi\epsilon)$.  The quantitative phase estimate
\cite[Lemma~2.9]{BayerDeformation2019} then makes the slicing distance
less than $\epsilon$ as long as it is less than $1/4$.  It is zero at the start;
continuity along the path and $\epsilon<1/8$ prevent a first exit from
that $1/4$ neighborhood.  This proves the asserted strict distance bound.
The charge bound is strict, so for each fixed $t$ it also holds on an open
neighborhood in \emph{all} ambient numerical charges.  This is the space
on which we deform, before applying restriction.
\end{proof}

\begin{proposition}\label{prop:geometric-prescribed-charges}
For either factor in \eqref{eq:two-character-factors}, the choices in
Lemma~\ref{lem:uniform-ambient-deformation} can be made so that every
$B\in V,H\in W,t\ge T$ gives a locally finite geometric stability
condition on $X$, with charge \eqref{eq:exact-restricted-charge} and full
numerical support.  Every point sheaf is stable of phase $1$.
For $\Theta=1$, the construction also gives compatible geometric slicings
on every smooth section $S_h\in|L_h|$, with Harder--Narasimhan phases
computed by pushforward to $Y$.
\end{proposition}

\begin{proof}
Set $A=\OO_Y(1,\ldots,1)$. Choose a locally free resolution
$P^\bullet\to i_*\OO_X$ with $P^0=\OO_Y$ and with each $P^{-k}$ a
finite sum of line bundles $A^{-b}$ for $k\ge1$. Such a resolution is
obtained by repeatedly generating the successive coherent kernels after
sufficiently positive twists. Fix $N>\dim Y+1$ and truncate after the
terms in degrees $0,\ldots,-(N-1)$. The remaining kernel is shifted by
$[N]$, so the cofiltration remainder lies in $D^{\le-N}(Y)$.

For a defining ambient linear form of a smooth section $S_h$, first form
the cone of the untruncated map
$P^\bullet(-\xi_h)\to P^\bullet$. It resolves
$(i\circ j_h)_*\OO_{S_h}$, and its term in degree $-k$ is
\[
 P^{-k}\oplus P^{1-k}(-\xi_h),\qquad P^1=0.
\]
Truncate each cone resolution after its terms in degrees
$0,\ldots,-(N-1)$. Its remaining kernel is again shifted by $[N]$, so
its cofiltration remainder also lies in $D^{\le-N}(Y)$.
The line bundles in all these terms are independent of the chosen smooth
section. Include all their summands, in particular the mixed twists
$A^{-b}\otimes\OO_Y(-\xi_h)$, together with $A$ and $\OO_Y(\xi_h)$,
in one finite list. This list is fixed before the parameters $B,H,t$ and
the sections are chosen.

Let $\mathcal R$ be a lift from
Lemma~\ref{lem:uniform-ambient-deformation}.  Tensoring is an isometry for
slicing distance, so \eqref{eq:ambient-twist-distance} gives, after
increasing the single threshold $T$,
\begin{equation}\label{eq:deformed-twist-distance}
 d(\mathcal R,\mathcal R\otimes L)<1
\end{equation}
for every line bundle $L$ in the finite list, uniformly in the claimed
sector.  The triangle inequality bounds the left side by
$2\epsilon+O(t^{-1})$.  The same strict inequalities hold in a full open
neighborhood of each lifted charge.

In particular tensoring by $A^{-1}[1]$ has the Bayer property.
\cite[Proposition~3.27]{LiEtAl2026} makes every point sheaf on $Y$
semistable throughout that open neighborhood.  Its numerical class is
primitive, because $\chi(\OO_Y,\OO_y)=1$.
Lemma~\ref{lem:primitive-point} upgrades semistability to stability.
At the charge \eqref{eq:deformed-ambient-charge} its value is $-1$, and
closeness to $\mathcal R_t$ fixes its phase to be $1$.

We can now apply Lemma~\ref{lem:geometric-range} on $Y$.
It supplies \eqref{eq:restriction-hypotheses} with
$n_1=-\dim Y$ and $n_2=0$.  Each resolution term $L[k]$, $k\ge1$, has
the Bayer property: \eqref{eq:deformed-twist-distance} says that tensoring
by $L$ lowers no extremal phase by as much as one, and the shift $[k]$
restores the required lower bound.  The initial $\OO_Y$ term is the
identity.  The chosen $N$ exceeds the truncation bound.  Thus the
finite-map criterion applies to the closed immersion $i$ and, when needed,
to every $i\circ j_h$.  It supplies slicings with the indicated charges
and Harder--Narasimhan filtrations.  Definition
\eqref{eq:induced-slicing} and uniqueness of those filtrations show that
pushforward computes their extremal phases.

Pushforward is exact for these hearts and detects zero objects.  A proper
destabilizing subobject of a point would therefore give one for the stable
point on $Y$.  The induced points are stable of phase $1$.
If $E$ is semistable on $X$, the ambient support inequality and the
injectivity in Lemma~\ref{lem:full-numerical-embedding} give
\[
 \|[E]\|\le C\|i_*[E]\|\le C'|U_t^{B,H}(i_*E)|.
\]
This is support on the full Euler numerical group of $X$.  For each
auxiliary surface the ambient lattice through pushforward suffices; no
injectivity assertion for its full numerical group is needed.
Finally the support bound on a finite-rank lattice gives local finiteness:
in a sufficiently short phase interval, the projection of a nonzero
semistable charge onto its middle ray has a fixed positive lower bound.
Additivity of this projection bounds the length of any chain of strict
subobjects or quotients of a fixed object in that interval.  The induced
slicings are thus locally finite, as asserted.
\end{proof}

Apply the construction to both choices of $\Theta$, intersect the two
neighborhoods $V$ and $W$, and take the larger of their thresholds.  The
same real sector and one threshold therefore work for both charges.
The proposition completes the prescribed Todd-charge assertion of
Theorem~\ref{thm:geometric-sector}.  It also supplies the exact ordinary
charge and full numerical support.  Identifying its heart requires the
additional argument in the next section.  Only $K_X\simeq\OO_X$ has been
used; no vanishing of $H^1(X,\OO_X)$ or $H^2(X,\OO_X)$ is required.
